\documentclass[11pt,a4paper]{article}
\usepackage[utf8]{inputenc}
\usepackage[T1]{fontenc}
\usepackage[brazil]{babel}
\usepackage[a4paper,margin=2.3cm]{geometry}
\usepackage{enumitem}
\usepackage{hyperref}
\usepackage{xcolor}
\setlength{\parindent}{0pt}
\setlength{\parskip}{5pt}
\sloppy
\setlength{\emergencystretch}{3em}
\hypersetup{colorlinks=true,urlcolor=blue,linkcolor=black}
\newcommand{\cabecalho}[2]{\begin{center}{\small \textbf{Sprint 2 -- #1}}\\[1mm]\textit{Inventário de Ativos e Vulnerabilidades}\\Matheus Moreli Durão Silva -- 12621CBS217\\Cibersegurança -- Bacharelado -- Noturno\end{center}\hrule\vspace{4mm}}
\newcommand{\bloco}[1]{\par\vspace{2mm}\textbf{#1}\par}
\newcommand{\codigo}[1]{\texttt{#1}}

\begin{document}\cabecalho{S2\_18 -- Testes e Correções}{}
\bloco{Ferramenta usada} \codigo{unittest} já vem com Python; não exige instalação externa. Cada teste monta sua própria base no método \codigo{setUp}, então o resultado não depende do JSON real, de internet ou de dados sensíveis.
\bloco{O que é testado} Cadastro e busca por ID; duplicidade; filtro sem modificação; atualização preservando ID e criando histórico; exclusão cancelada e confirmada; vulnerabilidade válida; CVE/CVSS inválidos; ida e volta do JSON em arquivo temporário.
\bloco{Como executar} Na pasta do projeto, usar \codigo{python -m unittest -v test\_sprint2.py}. O \codigo{-v} mostra o nome de cada teste; se algo falhar, o \codigo{assert} informa o valor esperado e o encontrado. A contagem verificada em 13/10/2026 foi 36 testes nesse arquivo, todos com \codigo{OK}.
\bloco{Correção sem regressão} Após uma mudança, executar a suíte inteira. Isso verifica que corrigir uma regra, como CVSS, não quebrou o cadastro de ativos ou a persistência.
\bloco{Cobertura do CRUD} Cadastro: válido, ID e nome duplicados. Consulta: ID presente/ausente e filtro sem alteração. Atualização: válida e inválida sem alteração parcial. Exclusão: cancelada e confirmada, inclusive remoção de vulnerabilidades e dependências. Cada operação tem ao menos dois cenários automatizados. O histórico de falha/correção deve usar saídas reais de execução; não há registro anterior de um teste vermelho neste entregável.
\end{document}
